\section{Wrap-up}
\label{section:tailoring:wrap-up}

The points laid out here are to account for two main scenarios: (1)
``low-hanging fruit" in performance, e.g., gains may be found immediately by
simply shifting input data and increasing I/O bandwidth; (2) ensuring that
there is greater stability in our analysis before progressing. This second
point is subtle, but whenever benchmarking we are trying to get machines to
behave in as controlled and predictable a way as possible. In the case of
thread and rank placement, we may see variable performance due to low thread
affinity or inconsistent memory locations, making deeper-level analyses harder
to rely on. 

This is not intended to be fully exhaustive: there are many dials we can tune
for a code, some can be impactful whilst others will be negligible. We have
listed here some initial variables to consider for tailoring a code that we
recommend testing in conjunction with the high-level metrics. For example, if
we try out hardware-specific compiler optimisations, we advise recomputing the
core-level metric to see if there is a significant change. In fact, we may see
shifts in other metrics too. The key is to test these features and document what we
see, using the metrics as a guide to track changes in performance.
